\documentclass[11pt]{article}

\usepackage{amsmath,amssymb,amsthm}
\usepackage{graphicx}
\usepackage{hyperref}
\usepackage{booktabs}
\usepackage{enumitem}

\newtheorem{theorem}{Theorem}
\newtheorem{conjecture}{Conjecture}
\newtheorem{proposition}{Proposition}
\newtheorem{lemma}{Lemma}
\newtheorem{definition}{Definition}
\newtheorem{remark}{Remark}

\title{Maximal Volume Ideal Polyhedra and the Arithmetic Angle Phenomenon}

\author{Igor Rivin}

\date{\today}

\begin{document}

\maketitle

\begin{abstract}
We present a software suite for the analysis and optimization of ideal convex polyhedra in hyperbolic 3-space $\mathbb{H}^3$. Using Rivin's variational characterization of ideal polyhedra, we develop efficient algorithms for checking combinatorial realizability and finding volume-maximizing configurations. Our systematic computational study reveals two striking phenomena: (1) maximal volume ideal polyhedra consistently exhibit dihedral angles that are rational multiples of $\pi$---a property with no obvious explanation from the optimization formulation; and (2) the distribution of volumes for random configurations is well-approximated by a Beta distribution, with mean normalized volume converging to approximately $\ln 2 \approx 0.69$ as the vertex count increases. We provide complete data for small vertex counts, including vertex positions, triangulations, and verified rational angle structures. Extending to $\mathbb{H}^4$ and $\mathbb{H}^6$, we discover a \emph{combinatorial invariance} phenomenon: the $f$-vector of the volume-maximizing ideal polytope with $N$ vertices on $S^{d-1}$ is the same whether volume is measured in the Euclidean or hyperbolic sense, holding exactly for $N \lesssim 5d$ (tested up to $d=3$, $N \leq 50$; $d=4$, $N \leq 40$; $d=5$, $N \leq 20$; $d = 6$, $N \leq 30$).  At larger~$N$, small deviations appear ($\leq 3$ facets), with the hyperbolic optimizer consistently finding fewer facets.  CMA-ES (Covariance Matrix Adaptation) provides Euclidean warm starts that are within $0.0$--$0.4\%$ of the hyperbolic optimum, confirming the near-equivalence of the two objectives near their shared maximum. An interactive implementation is publicly available.
\end{abstract}

\smallskip\noindent
\textbf{Keywords:} ideal polyhedra, hyperbolic volume, Lobachevsky function,
dilogarithm, combinatorial invariance, Frobenius groups, vertex-transitive
polytopes, CMA-ES, Schl\"afli formula

\smallskip\noindent
\textbf{MSC 2020:} 52B10 (primary), 51M25, 52A40, 52B15, 20B25, 11R18

\tableofcontents

% ============================================================================
\section{Introduction}
% ============================================================================

Ideal polyhedra in hyperbolic 3-space $\mathbb{H}^3$ occupy a central position in low-dimensional topology and geometry. Their study connects to the geometry of hyperbolic 3-manifolds, Teichmüller theory, and the representation theory of surface groups. The volume of hyperbolic 3-manifolds, a topological invariant, can often be computed by decomposing the manifold into ideal tetrahedra---the simplest ideal polyhedra.

A natural question arises: among all ideal polyhedra with a given number of vertices, which has the largest volume? This optimization problem, while analytically straightforward (it reduces to maximizing a concave function over a convex polytope), reveals surprising structure. The maximal configurations exhibit a form of ``arithmeticity''---their dihedral angles are consistently rational multiples of $\pi$. Moreover, the distribution of volumes for random configurations follows a Beta distribution, with parameters scaling linearly in the vertex count and mean normalized volume converging to approximately $\ln 2$.

In this paper, we present a comprehensive software suite for:
\begin{enumerate}[label=(\roman*)]
    \item Checking whether a given triangulation is realizable as the boundary of an ideal convex polyhedron
    \item Finding optimal (volume-maximizing) geometric realizations of realizable triangulations
    \item Computing the Fuchsian group associated with the surface of the polyhedron
    \item Visualizing polyhedra in multiple models of hyperbolic space
\end{enumerate}

Our computational experiments reveal a surprising phenomenon: when optimizing for maximal volume, the resulting polyhedra consistently have dihedral angles that are \emph{exact rational multiples of $\pi$}. This arithmetic structure emerges from a purely analytic optimization problem with no built-in number-theoretic constraints.

% ============================================================================
\section{Background and Notation}
% ============================================================================

\subsection{Ideal Polyhedra in $\mathbb{H}^3$}

An \emph{ideal polyhedron} in hyperbolic 3-space $\mathbb{H}^3$ is a convex polyhedron whose vertices lie on the boundary at infinity $\partial_\infty \mathbb{H}^3 \cong \mathbb{CP}^1$. Despite having vertices ``at infinity,'' such polyhedra have finite volume---a consequence of the exponential shrinking of hyperbolic space toward the boundary.

The combinatorial structure of an ideal polyhedron is encoded by a triangulation of the 2-sphere: the \emph{Delaunay triangulation} of the ideal vertices as seen from the center of mass of the polyhedron. Conversely, not every triangulation of the sphere arises from an ideal polyhedron---realizability imposes constraints on the possible angle structures, as characterized by Rivin's theorem below.

We work primarily in the upper half-space model, where $\mathbb{H}^3 = \{(x,y,z) \in \mathbb{R}^3 : z > 0\}$ with metric $ds^2 = (dx^2 + dy^2 + dz^2)/z^2$. The boundary $\partial_\infty \mathbb{H}^3$ is the extended complex plane $\mathbb{C} \cup \{\infty\}$, and ideal vertices are specified as points in $\mathbb{CP}^1$.

\subsection{Rivin's Variational Characterization}

The foundational result for our approach is Rivin's characterization of ideal polyhedra:

\begin{theorem}[Rivin \cite{Rivin1994,RivinCombOpt}]
A triangulation $\mathcal{T}$ of the sphere is realizable as the boundary of a convex ideal polyhedron in $\mathbb{H}^3$ if and only if there exist positive real numbers $\theta_c > 0$ assigned to each corner $c$ of each triangle such that:
\begin{enumerate}[label=(\roman*)]
    \item \textbf{Triangle constraint:} For each triangle $T$ with corners $c_1, c_2, c_3$:
    \[
    \theta_{c_1} + \theta_{c_2} + \theta_{c_3} = \pi
    \]
    \item \textbf{Vertex constraint:} For each vertex $v$, summing over all corners at $v$:
    \[
    \sum_{c \text{ at } v} \theta_c = 2\pi
    \]
    \item \textbf{Edge constraint:} For each edge $e$, let $\alpha_e$ and $\beta_e$ be the angles opposite to $e$ in the two triangles sharing $e$. Then:
    \[
    \alpha_e + \beta_e < \pi
    \]
\end{enumerate}
\end{theorem}

The angles $\theta_c$ are the \emph{dihedral angles} of the ideal polyhedron---the angle between adjacent faces along their common edge, measured inside the polyhedron. The key insight is that these constraints are all linear (or strict linear inequalities), so realizability can be checked via linear programming, and volume optimization becomes a concave maximization problem over a convex polytope.

\subsection{The Lobachevsky Function and Hyperbolic Volume}

The hyperbolic volume of an ideal tetrahedron with dihedral angles $\alpha, \beta, \gamma$ (opposite pairs equal) is given by:
\[
V(\alpha, \beta, \gamma) = \Lambda(\alpha) + \Lambda(\beta) + \Lambda(\gamma)
\]
where $\Lambda(\theta)$ is the Lobachevsky function:
\[
\Lambda(\theta) = -\int_0^\theta \log |2\sin t| \, dt
\]

The Lobachevsky function satisfies $\Lambda(\pi - \theta) = -\Lambda(\theta)$ and has maximum value $\Lambda(\pi/3) \approx 0.5074$ (Catalan's constant divided by 2). The function is strictly concave on $(0, \pi)$, which ensures that the volume function over the Rivin polytope has a unique maximum. The regular ideal tetrahedron, with all dihedral angles equal to $\pi/3$, achieves volume $V_4 = 3\Lambda(\pi/3) \approx 1.0149$.

\subsection{The Penner-Rivin Algorithm}

Given a realized ideal polyhedron with its dihedral angles, we can compute the associated Fuchsian group $\Gamma < \mathrm{PSL}(2,\mathbb{R})$ using the Penner-Rivin algorithm \cite{RivinTriangulations}. The surface of the ideal polyhedron inherits a complete hyperbolic metric, making it a hyperbolic surface $S$ with cusps at the ideal vertices. The fundamental group $\pi_1(S)$ acts on the hyperbolic plane via a discrete faithful representation---the \emph{holonomy representation}.

The algorithm proceeds as follows:
\begin{enumerate}[label=(\roman*)]
    \item \textbf{Decorate the triangulation:} Assign to each edge $e$ the complex number
    \[
    z_e = \exp(i(\alpha_e + \beta_e))
    \]
    where $\alpha_e, \beta_e$ are the angles opposite $e$ in the adjacent triangles. This is the \emph{shape parameter} of the edge.

    \item \textbf{Compute edge matrices:} For each oriented edge from vertex $u$ to vertex $v$, construct a $2 \times 2$ matrix encoding the hyperbolic translation along the edge, parameterized by the shape parameters of the surrounding edges.

    \item \textbf{Compose around loops:} For each generator of $\pi_1(S)$, represented as a path crossing edges of the triangulation, multiply the corresponding edge matrices to obtain the holonomy matrix in $\mathrm{PSL}(2,\mathbb{R})$.
\end{enumerate}

The resulting matrices generate the Fuchsian group $\Gamma$. When all shape parameters have modulus 1 (equivalently, when all $\alpha_e + \beta_e$ are real), the group is \emph{arithmetic} if certain additional number-theoretic conditions hold. Our observation that maximal volume configurations have rational angles $p\pi/q$ suggests a connection to arithmetic Fuchsian groups, since in these cases $z_e = \exp(ip\pi/q)$ is a root of unity.

% ============================================================================
\section{Methods}
\label{sec:methods}
% ============================================================================

We describe the computational methods in rough order of generality,
beginning with the $d = 3$ angle-based approach (which is exact
and special to $\mathbb{H}^3$) and proceeding to the vertex-based
optimization pipeline used in dimensions $d = 4, 5, 6$.

\subsection{Dimension 3: Angle-Based Optimization}

\subsubsection{Realizability testing}

Given a triangulation $\mathcal{T}$ with $n$ vertices, $f$ faces,
and $e$ edges, realizability as an ideal polyhedron is equivalent
to feasibility of a linear program.  Introduce variables $\theta_c$
for each corner $c$ (there are $3f$ corners).  The constraints from
Theorem~1 become:
\begin{align*}
\theta_{c_1} + \theta_{c_2} + \theta_{c_3} &= \pi
  && \text{for each triangle } T, \\
\sum_{c \text{ at } v} \theta_c &= 2\pi
  && \text{for each vertex } v, \\
\alpha_e + \beta_e &< \pi
  && \text{for each edge } e, \\
\theta_c &> 0
  && \text{for each corner } c.
\end{align*}
Strict inequalities are implemented with a tolerance
$\epsilon = 10^{-6}$.

\subsubsection{Volume optimization in angle coordinates}

The volume is a concave function of the dihedral angles:
\[
V(\boldsymbol{\theta})
= \sum_T \bigl(\Lambda(\theta_{c_1})
  + \Lambda(\theta_{c_2})
  + \Lambda(\theta_{c_3})\bigr),
\]
where $\Lambda$ is the Lobachevsky function.  Concavity over the
convex Rivin polytope guarantees a unique global maximum, which we
find by L-BFGS-B from multiple random starting points.

\subsubsection{Direct vertex optimization via Bloch--Wigner}

For larger vertex counts ($n \geq 20$), we bypass the angle
parameterization and optimize directly over vertex configurations
on~$S^2$.  Each ideal tetrahedron $[\infty, z_1, z_2, z_3]$ has
volume $|D(z)|$ where $D$ is the Bloch--Wigner dilogarithm and
$z = (z_1 - z_3)/(z_2 - z_3)$ is the cross-ratio.  The
Bloch--Wigner function is evaluated via
$D(z) = \mathrm{Im}(\mathrm{Li}_2(z)) + \arg(1-z)\log|z|$
using \texttt{scipy.special.spence}, giving $\sim 0.1\,\text{ms}$
per volume evaluation.

CMA-ES~\cite{HansenOstermeier2001} is applied directly to the
$\mathbb{H}^3$-volume objective (not the Euclidean proxy), with
20 restarts and $500n$ function evaluations per restart.  For
$n \geq 100$, sep-CMA-ES (diagonal covariance approximation) is
used, reducing the per-generation cost from $O(n^2)$ to~$O(n)$.

\subsubsection{Rational angle detection}

We detect rational multiples of $\pi$ via continued fraction
expansion of $\theta/\pi$.  If some convergent $p_k/q_k$
satisfies $|\theta/\pi - p_k/q_k| < 10^{-10}$ with $q_k \leq 100$,
we report $\theta = p_k\pi/q_k$.  All maximal configurations
tested pass this test at machine precision.

\subsection{Dimensions $d \geq 4$: Vertex-Based Optimization}

In $d \geq 4$, the Rivin angle parameterization is unavailable:
the configuration space is the set of $N$ points on $S^{d-1}$
(ideal vertices), and the volume must be computed from vertex
coordinates.  Our pipeline has three stages.

\subsubsection{Stage 1: Euclidean warm start via CMA-ES}

We maximize the Euclidean convex hull volume of $N$ points
on $S^{d-1}$ using the Covariance Matrix Adaptation Evolution
Strategy (CMA-ES)~\cite{HansenOstermeier2001}.  CMA-ES adapts a
multivariate Gaussian distribution to the promising region of
parameter space, without requiring gradients.  The objective
function is $-\mathrm{Vol}(\mathrm{conv}(x_1,\ldots,x_N))$,
evaluated via \texttt{scipy.spatial.ConvexHull} (which calls
Qhull's C implementation).

The Cartesian coordinates of the $N$ points are the $Nd$ free
parameters; after each CMA-ES step, vertices are projected back to
$S^{d-1}$.  Typical budget: $500N$ function evaluations with initial
step size $\sigma_0 = 0.3$, seeded from the best of 5--10 random
configurations.  CMA-ES dramatically outperforms random restarts:
in $d = 4$, the Euclidean volume of the CMA-ES solution exceeds
that of 50 random restarts by 40--230\%, and is already within
$0.0$--$0.4\%$ of the hyperbolic optimum.

\subsubsection{Stage 2: Hyperbolic polish}

The CMA-ES Euclidean solution is used as a warm start for
hyperbolic volume maximization.  The approach depends on
the dimension.

\paragraph{$d = 4$: Haagerup--Munkholm formula + Powell.}
The ideal 4-simplex volume is given by the Haagerup--Munkholm
angle formula~\cite{HM}, involving Euclidean dihedral angles,
circumradii, and signed distances---all closed-form from the
vertex coordinates.  The polytope volume is the sum over a
Delaunay triangulation into ideal 4-simplices.  Since each volume
evaluation takes $\sim 2\,\text{ms}$, we use Powell's derivative-free
method for the hyperbolic polish.

\paragraph{$d = 5$: Schl\"afli gradient + LBFGS.}
In odd dimensions there is no closed-form volume formula.
However, \emph{for optimization the volume itself is not needed}.
The Schl\"afli differential formula
\[
  dV_5 = -\tfrac{1}{4}\sum_{F} V_3(F)\,d\theta_F
\]
gives the exact gradient in closed form: the codimension-2 faces
$F$ of an ideal 5-simplex are ideal tetrahedra in a totally
geodesic $\mathbb{H}^3$, with volumes given by the
Milnor--Lobachevsky formula $V_3 = \Lambda(\alpha) + \Lambda(\beta)
+ \Lambda(\gamma)$, and the dihedral angles $\theta_F$ come from
the angle Gram matrix.  We build a ``surrogate'' function
$S = -\frac{1}{4}\sum_F V_3(F) \cdot \theta_F$ whose gradient
equals the Schl\"afli gradient exactly (treating $V_3(F)$ as
constants), and optimize $S$ with torch LBFGS.  The actual
volume is computed once at the final configuration via the
Schl\"afli path integral.

\paragraph{$d = 6$: Hopf--Murakami formula + torch LBFGS.}
For even dimensions, the Poincar\'e--Schl\"afli--Hopf reduction
gives a closed-form volume for each ideal 6-simplex:
\[
  V_6 = \frac{8\pi^3}{15}\Bigl(\frac{17}{4} - w_4 + \frac{1}{2}w_2\Bigr),
\]
where $w_4$ involves dihedral angles and $w_2$ involves $S^3$
link volumes computed by Murakami's dilogarithm formula.  We
implement this in PyTorch with a custom \texttt{autograd.Function}
for the dilogarithm $\mathrm{Li}_2$: a pure-torch forward pass
using Bernoulli-series evaluation with argument reduction
(inversion + reflection identities map all arguments to the
convergence region), and an analytical backward pass
$d\mathrm{Li}_2/dz = -\log(1-z)/z$.  All operations are
vectorized over simplices and batched over the $\binom{7}{4} = 35$
Murakami calls per 6-simplex, yielding a single autograd node per
simplex.  On an NVIDIA GH200 GPU, the full forward + backward
pass for $N = 80$ (3175 simplices, 889\,000 $\mathrm{Li}_2$
evaluations) takes $\sim 16\,\text{ms}$.

\subsubsection{Stage 3: Triangulation management}

In all dimensions $d \geq 4$, the ideal polytope is decomposed
into simplices via stereographic projection from one vertex
(the ``pole'') to $\mathbb{R}^{d-1}$, followed by Delaunay
triangulation of the projected points.  Coning from the pole
vertex produces $(d+1)$-vertex simplices suitable for the volume
formula.  The triangulation is recomputed periodically during
LBFGS optimization (every 5 steps) to accommodate changes in
the combinatorial type; the cost of a full Delaunay recomputation
(via Qhull) is negligible compared to the volume evaluation.

\subsubsection{Structured initialization: a small-$N$ phenomenon}
\label{sec:structured-seeds}

In addition to the random Stage~1 starts, we maintain a library
of \emph{structured} initial configurations---free sums of regular
polygons in mutually orthogonal $2$-planes
(\S\ref{sec:joins}), the cross-polytope, the regular simplex and
its joins, and (for $d = 4$) the $24$-cell.  Each natural object
has a fixed vertex count; at $N$ slightly above that count we
augment with a few random points, but only up to the next natural
object's size.  We run the hyperbolic polish from each member of
this library and take the best.

Empirically, structured seeds matter substantially at small $N$
and fade---then invert---as $N$ grows.  At $N \leq 2d + 2$ the
volume-optimal polytope is one of the structured library members
verbatim (see Table in \S\ref{sec:joins}: $N = d{+}2$ is
$\Delta_{d/2}{*}\Delta_{d/2}$, $N = 2d$ is the cross-polytope, etc.),
and structured warm starts converge to the optimum on the first
polish, where random restarts often miss it entirely.  Past the
library-exhaustion point---empirically near $N \approx 2d^2$,
where the largest natural object in our library
($24$-cell at $d = 4$, $C_4^3$ cross-polytope at $d = 6$) has
been outgrown---the relationship inverts.  In a controlled
$1{,}000$-restart comparison at $d = 4$, the best structured-seed
final volume is $97\%$ of the best random-restart final volume at
$N = 30$, $94\%$ at $N = 50$, and $90\%$ at $N = 80$; at
$d = 6, N = 35$ the ratio drops to $77\%$.  The structured-seed
advantage is genuinely a small-$N$, low-dimensional artefact.

The clean interpretation is Erd\H{o}s--R\'enyi-style.  At small $N$
the moduli space of inscribed simplicial polytopes is small enough
that the volume-optimal stratum is forced to contain an
algebraically distinguished configuration; structured seeds find
this stratum by construction, while random restarts have to hit it
by chance.  As $N$ grows, the number of combinatorial types grows
at least at Catalan rate, the global optimum's automorphism group
collapses to triviality (Tables~\ref{tab:comb_d4}
and~\ref{tab:comb_H6}), and the volume-optimal stratum is no
longer algebraically distinguished---random restarts now sample
relevant strata at a rate proportional to their parameter-space
measure, which is the right rate for high-$N$.  The structured-seed
advantage tracks the ``law of small numbers'': there are only
finitely many natural algebraic objects, and once $N$ exceeds them
the optimization landscape ceases to know about them.

\subsection{Symmetry Analysis}

The geometric automorphism group of an optimal polytope is computed
as follows:
\begin{enumerate}[label=(\roman*)]
\item \textbf{Centering.}  The Douady--Earle conformal
  barycenter~\cite{DouadyEarle} of the ideal vertices is
  computed by fixed-point iteration (M\"obius transport of the
  centroid), then moved to the origin of the Poincar\'e ball by
  the corresponding M\"obius isometry.  After centering, every
  geometric symmetry of the configuration is an element of
  $\mathrm{O}(d)$.
\item \textbf{Distance coloring.}  The chordal distances between
  centered vertices are clustered into equivalence classes
  (tolerance $10^{-4}$--$10^{-6}$), yielding a coloring of
  the complete graph $K_N$.
\item \textbf{Automorphism computation.}  The colored graph is
  encoded as a vertex-colored graph (via the standard
  gadget-vertex trick) and fed to
  \textsc{nauty}~\cite{McKayPiperno} via \texttt{pynauty}.
  The output is the automorphism group order, generators as
  permutations, and vertex orbits.
\end{enumerate}

\subsection{Hessian Analysis}

To classify critical points as local maxima, saddle points, or
degenerate, we compute the Hessian of the volume functional
restricted to the tangent space of $(S^{d-1})^N$.  At each vertex
$v_i$, a $(d-1)$-dimensional orthonormal tangent basis is
constructed by Gram--Schmidt projection of the standard basis
against $v_i$.  The Hessian is assembled by finite differencing
of the gradient in each tangent direction (for $d = 5, 6$, the
gradient itself is computed analytically via the Schl\"afli or
Hopf formula).  The resulting $(d-1)N \times (d-1)N$ symmetric
matrix is diagonalized; the signature (numbers of positive,
negative, and zero eigenvalues) determines the critical point type.
The expected nullity is $\dim\mathrm{SO}(d) = d(d-1)/2$ from
rotational invariance; any additional zero eigenvalues indicate
a volume-constant ridge in the configuration space.

% ============================================================================
\section{Results}
% ============================================================================

\subsection{Maximal Volume Configurations}

We computed maximal volume ideal polyhedra for $n = 4, 5, \ldots, 12$ vertices using 100 independent optimization trials per vertex count, and extended the results to much larger~$n$ using CMA-ES with the Bloch--Wigner
dilogarithm evaluator ($\sim 0.1\,\text{ms}$ per volume computation).
Table~\ref{tab:maximal-volumes} summarizes our findings.

\begin{table}[htbp]
\centering
\caption{Maximal volume ideal polyhedra with $n$ vertices.
$V_4 \approx 1.0149$ is the volume of the regular ideal tetrahedron.
For $n \leq 12$, angle denominators are shown; for larger~$n$ the
volume was maximized directly via CMA-ES on the Bloch--Wigner
functional over vertex configurations on~$S^2$.}
\label{tab:maximal-volumes}
\begin{tabular}{ccccl}
\toprule
$n$ & Volume & Faces & $V/V_4$ & Notes \\
\midrule
4 & 1.014942 & 4 & 1.0000 & $q = 3$ \\
5 & 2.029883 & 6 & 2.0000 & $q = 3$ \\
6 & 3.663862 & 8 & 3.6099 & $q = 2$ \\
7 & 4.986773 & 10 & 4.9134 & $q = 5$ \\
8 & 6.488469 & 12 & 6.3929 & $q = 11$ \\
9 & 8.162538 & 14 & 8.0424 & mixed \\
10 & 9.839315 & 16 & 9.6946 & $q = 14$ \\
11 & 11.449290 & 18 & 11.2810 & $q = 7, 14$ \\
12 & 13.529628 & 20 & 13.3308 & $q = 5$ \\
\midrule
20 & 27.699 & 36 & 27.29 & \\
30 & 46.486 & 56 & 45.80 & \\
50 & 84.971 & 96 & 83.72 & \\
70 & 123.893 & 136 & 122.07 & \\
100 & 182.671 & 196 & 180.02 & sep-CMA-ES \\
150 & 279.727 & --- & 275.61 & sep-CMA-ES \\
200 & 377.729 & --- & 372.17 & sep-CMA-ES \\
300 & 573.429 & --- & 565.00 & sep-CMA-ES \\
\bottomrule
\end{tabular}
\end{table}

As an independent check, a $300$-restart vertex-space search
(random initialization on $S^2$, CMA-ES on the hyperbolic objective,
scipy Powell polish) reproduces the $n = 11$ and $n = 12$ records
to within $10^{-4}$ and $5 \cdot 10^{-4}$ respectively, confirming
that the angle-parameterization maxima reported above are not
artifacts of a small search budget.

For $n \geq 100$, we use sep-CMA-ES (CMA-ES with diagonal
covariance), which scales as $O(n)$ per generation rather than
$O(n^2)$ for the full covariance update.  Covariance spectrum
analysis reveals that the volume landscape on $\mathrm{Conf}_n(S^2)$
has only 2 dominant eigenvalues, with the remaining $3n - 2$
directions being nearly isotropic (condition number $\lesssim 20$).
This explains why the diagonal approximation loses less than $1\%$
in volume quality while being $2\times$ faster.

The per-vertex volume contribution $V_n/n$ increases monotonically:
$1.82$ at $n = 100$, rising to $1.91$ at $n = 300$.

\begin{remark}[Comparison with Platonic solids]
For $n=6$, the ``obvious'' candidate is the ideal right-angled octahedron, with all dihedral angles equal to $\pi/2$. Indeed, our optimization confirms this: the maximal configuration has angle denominator $q=2$, meaning all angles are multiples of $\pi/2$. The right-angled octahedron is optimal among 6-vertex ideal polyhedra.

For $n=8$, the natural candidate is the ideal cube (with vertices at the 8 corners of a cube inscribed in the sphere). The ideal cube has $V/V_4 = 5$, giving volume $V_{\text{cube}} \approx 5.07$. However, our optimization finds a configuration with $V \approx 6.49$, exceeding the cube by approximately 28\%. The optimal 8-vertex polyhedron is \emph{not} combinatorially a cube---it has a different triangulation with the surprising angle denominator $q=11$.

For $n=12$, the natural candidate is the ideal regular icosahedron. The icosahedron has $V/V_4 = 10$, giving $V_{\text{icosahedron}} \approx 10.15$. Our optimization finds $V \approx 13.53$, exceeding the icosahedron by approximately 33\%. As with the cube, the optimal 12-vertex polyhedron has a different combinatorial structure. Remarkably, all dihedral angles are multiples of $\pi/5$: specifically $2\pi/5$, $3\pi/5$, $4\pi/5$, and $\pi$ (a flat edge).
\end{remark}

\begin{remark}[Symmetry of optimal configurations]
The automorphism groups of the optimal triangulations show a clear pattern: high symmetry for small $n$, rapidly decreasing thereafter.
\begin{center}
\begin{tabular}{ccc}
$n$ & $|\mathrm{Aut}|$ & Structure \\
\hline
4 & 6 & $S_3$ (tetrahedron) \\
5 & 6 & $D_3$ (bipyramid) \\
6 & 8 & $D_4$ (octahedron) \\
7 & 4 & $\mathbb{Z}_2 \times \mathbb{Z}_2$ \\
8--12 & 2--4 & $\mathbb{Z}_2$ or $\mathbb{Z}_2 \times \mathbb{Z}_2$
\end{tabular}
\end{center}
Volume maximization does not favor symmetric configurations---the optimal polyhedra for $n \geq 7$ have only a single reflection symmetry or none at all.
\end{remark}

\begin{figure}[htbp]
\centering
\includegraphics[width=\textwidth]{figures/triangulations_panel.png}
\caption{Planar Delaunay triangulations of the maximal volume ideal polyhedra. Each triangulation represents the combinatorial structure of the polyhedron projected onto the plane, with the point at infinity as the unbounded face.}
\label{fig:triangulations}
\end{figure}

\begin{figure}[htbp]
\centering
\includegraphics[width=\textwidth]{figures/poincare_panel.png}
\caption{Maximal volume ideal polyhedra visualized in the Poincar\'e ball model of $\mathbb{H}^3$. Vertices lie on the boundary sphere $\partial_\infty \mathbb{H}^3$, and the geodesic faces curve inward as characteristic of hyperbolic space.}
\label{fig:poincare}
\end{figure}

\subsection{The Rational Angle Phenomenon}

Our most striking finding is that volume-maximizing configurations \emph{within each combinatorial type} exhibit dihedral angles that are exact rational multiples of $\pi$. This holds not only for the globally maximal configurations shown in Table~\ref{tab:maximal-volumes}, but for every realizable triangulation we have tested.

\begin{conjecture}[Rational Angle Conjecture]
For every realizable triangulation $\mathcal{T}$ of the sphere, the volume-maximizing ideal polyhedron with combinatorial type $\mathcal{T}$ has all dihedral angles of the form $\frac{p\pi}{q}$ for integers $p, q$.
\end{conjecture}

This phenomenon is remarkable because:
\begin{itemize}
    \item The optimization has no built-in preference for rational angles
    \item The Lobachevsky function has no special behavior at rational multiples of $\pi$
    \item The constraint set (Rivin polytope) has no obvious arithmetic structure
    \item The phenomenon persists across \emph{all} combinatorial types, not just the globally optimal ones
\end{itemize}

For the globally maximal configurations shown in Table~\ref{tab:maximal-volumes}, we verified the rational angle property in all cases $n = 4, \ldots, 12$. The denominators range from small values like $q = 3$ (tetrahedron), $q = 2$ (octahedral $n=6$), and $q = 5$ ($n=7$ and $n=12$) to larger values like $q = 11$ for $n=8$ and $q = 14$ for $n = 10$ and $n=11$.

\subsection{Connection to Arithmetic Groups}

The prevalence of rational angles in maximal volume configurations suggests a deep connection to arithmetic hyperbolic geometry. When the dihedral angles are rational multiples of $\pi$, the shape parameters $z_e = \exp(i(\alpha_e + \beta_e))$ are roots of unity. This places strong constraints on the holonomy representation.

A Fuchsian group $\Gamma < \mathrm{PSL}(2,\mathbb{R})$ is \emph{arithmetic} if it is commensurable with the group of units in a quaternion algebra over a totally real number field. Arithmetic Fuchsian groups are characterized by having trace field $\mathbb{Q}(\mathrm{tr}(\gamma) : \gamma \in \Gamma)$ that is a number field, and satisfying additional ramification conditions.

When all shape parameters are roots of unity, the traces of holonomy matrices lie in cyclotomic fields. While this alone does not guarantee arithmeticity, it is highly suggestive. We conjecture that:

\begin{conjecture}[Arithmetic Holonomy]
For maximal volume ideal polyhedra, the holonomy representation $\rho: \pi_1(S) \to \mathrm{PSL}(2,\mathbb{R})$ has image commensurable with an arithmetic Fuchsian group.
\end{conjecture}

For the Platonic cases, the Fuchsian groups are known principal congruence subgroups:
\begin{itemize}
    \item $n=4$ (tetrahedron): $\Gamma(3)$
    \item $n=6$ (octahedron): $\Gamma(4)$
    \item Icosahedron: $\Gamma(5)$ --- but this is \emph{not} maximal for $n=12$
\end{itemize}
Identifying the Fuchsian groups corresponding to the maximal volume configurations for $n \geq 7$ remains an interesting open problem. The angle denominators (e.g., $q=11$ for $n=8$, $q=14$ for $n=10$) suggest these may be congruence subgroups of higher level.

\begin{remark}[Coxeter polytopes and branched covers]
A striking distinction emerges between the Platonic and non-Platonic cases. For $n=4$ and $n=6$, \emph{all} dihedral angles have the form $\pi/m$ for an integer $m$ (specifically $\pi/3$ and $\pi/2$ respectively). This means these ideal polyhedra are \emph{Coxeter polytopes}---fundamental domains for discrete reflection groups in $\mathbb{H}^3$.

For $n \geq 7$, however, the angles are $k\pi/q$ with $k > 1$. For instance, $n=7$ has angles $2\pi/5$, $3\pi/5$, $4\pi/5$, and $n=8$ has angles $4\pi/11$, $5\pi/11$, etc. These are \emph{not} Coxeter fundamental domains. The common denominator $q$ suggests these polyhedra may arise as fundamental domains for finite-index subgroups of reflection groups, or equivalently, that their universal covers are branched covers over some link complement. This connection to orbifold geometry and branched coverings merits further investigation.
\end{remark}

\subsection{Volume Distribution of Random Configurations}

To understand the landscape of ideal polyhedron volumes, we sampled 5000 random configurations for each vertex count $n \in \{5, 6, 7, 8, 10, 12\}$ by placing $n-3$ uniformly random points on the sphere (in addition to two fixed vertices and infinity) and computing the volume via the Delaunay triangulation.

Normalizing each volume by the maximal volume $V_{\max}(n)$ from Table~\ref{tab:maximal-volumes}, we find that the normalized volume distribution $V/V_{\max}$ is remarkably well-approximated by a Beta distribution:
\[
\frac{V}{V_{\max}} \sim \mathrm{Beta}(\alpha_n, \beta_n)
\]

Table~\ref{tab:beta-fits} summarizes the fitted parameters. The goodness of fit (KS test p-values $> 0.3$ for $n \geq 8$) confirms that the Beta distribution is an appropriate model.

\begin{table}[htbp]
\centering
\caption{Beta distribution fit parameters for normalized volume $V/V_{\max}$.}
\label{tab:beta-fits}
\begin{tabular}{ccccccc}
\toprule
$n$ & $\alpha$ & $\beta$ & Mean & Std & KS stat & p-value \\
\midrule
5  & 2.88  & 1.53  & 0.659 & 0.208 & 0.031 & 0.0001 \\
6  & 6.74  & 4.11  & 0.623 & 0.140 & 0.046 & 0.0000 \\
7  & 9.67  & 4.84  & 0.667 & 0.120 & 0.024 & 0.0053 \\
8  & 13.26 & 6.12  & 0.685 & 0.103 & 0.011 & 0.62 \\
10 & 23.02 & 10.05 & 0.696 & 0.079 & 0.012 & 0.48 \\
12 & 32.56 & 14.49 & 0.692 & 0.067 & 0.014 & 0.32 \\
\bottomrule
\end{tabular}
\end{table}

The parameters exhibit clear linear scaling with $n$:
\[
\alpha_n \approx 4.25n - 19.3, \qquad \beta_n \approx 1.78n - 7.4
\]
with the ratio $\alpha_n/\beta_n \approx 2.39$ remaining approximately constant. This implies that the mean normalized volume
\[
\mathbb{E}\left[\frac{V}{V_{\max}}\right] = \frac{\alpha_n}{\alpha_n + \beta_n} \approx 0.69 \approx \ln 2
\]
stabilizes as $n$ increases. The intriguing appearance of $\ln 2$ (or perhaps a related constant) as the asymptotic mean suggests a deeper connection to hyperbolic volume theory that we have not yet identified.

\begin{figure}[htbp]
\centering
\includegraphics[width=\textwidth]{figures/volume_distributions.png}
\caption{Volume distributions of random ideal polyhedra, normalized by $V_{\max}(n)$. Histograms (blue) show empirical distributions from 5000 samples; red curves show fitted Beta($\alpha$, $\beta$) distributions. The fit quality improves with $n$.}
\label{fig:volume-distributions}
\end{figure}

\begin{figure}[htbp]
\centering
\includegraphics[width=\textwidth]{figures/beta_parameters_scaling.png}
\caption{Scaling of Beta distribution parameters with vertex count $n$. Left: $\alpha$ grows approximately as $4.25n$. Center: $\beta$ grows as $1.78n$. Right: Mean normalized volume stabilizes near $0.69 \approx \ln 2$.}
\label{fig:beta-scaling}
\end{figure}

The concentration of the distribution (decreasing standard deviation as $n$ increases) indicates that for large $n$, most random configurations achieve volumes close to $69\%$ of the maximum---suggesting that volume optimization yields only modest improvements over random placement for large polyhedra.

% ============================================================================
\section{Visualizations}
% ============================================================================

Figure~\ref{fig:triangulations} shows the planar Delaunay triangulations for the maximal volume configurations with $n = 4, \ldots, 12$ vertices. These represent the combinatorial structure of the polyhedra, projected onto the plane via stereographic projection with one vertex (at infinity) as the unbounded outer face.

Figure~\ref{fig:poincare} shows the corresponding ideal polyhedra in the Poincar\'e ball model of $\mathbb{H}^3$. In this model, the ideal vertices lie on the boundary sphere, and the faces appear as geodesic surfaces that curve inward---a characteristic feature of negative curvature. The progression from the simple ideal tetrahedron ($n=4$) to more complex structures is clearly visible.

% ============================================================================
\section{Software}
% ============================================================================

The complete software suite is available at:
\begin{itemize}
    \item GitHub: \url{https://github.com/[repo]}
    \item Interactive demo: \url{https://huggingface.co/spaces/igriv/idealpolyhedra}
\end{itemize}

The implementation includes:
\begin{itemize}
    \item Realizability checking via Rivin's LP characterization
    \item Volume optimization via CMA-ES on the Euclidean surrogate
          and L-BFGS on the Bloch--Wigner / Murakami differentiable
          hyperbolic objective, with dimension-specific gradient paths
          for $d \in \{3,4,5,6\}$
    \item Fuchsian group computation via the Penner-Rivin algorithm
    \item 3D visualization in Klein, Poincaré ball, and sphere models
    \item Symmetry analysis using nauty
\end{itemize}

\paragraph{Development methodology.} Mathematical software has long suffered from ``bit rot''---the gradual degradation of code as dependencies evolve and environments change. This software was developed using an AI-assisted ``self-maintaining'' approach with Claude Code \cite{RivinDjinn}, which enables automated testing, dependency management, and iterative refinement. This methodology addresses the endemic problem of unmaintained research code by building in mechanisms for autonomous repair and adaptation.

% ============================================================================
% ============================================================================
\section{Higher Dimensions: Ideal Polytopes in $\mathbb{H}^4$}
\label{sec:H4}
% ============================================================================

\subsection{Setup}

An \emph{ideal polytope} in $\mathbb{H}^4$ has vertices on $S^3 = \partial_\infty \mathbb{H}^4$.
Given $N$ points on $S^3$, the ideal polytope is determined by the convex hull
of these points in the Klein model (the open unit ball in $\mathbb{R}^4$).
The 4-volume is computed by stereographic projection from one vertex to $\mathbb{R}^3$,
Delaunay triangulation, and the Haagerup--Munkholm angle formula for each ideal 4-simplex:
\[
\mathrm{Vol}_4[\infty, v_1, v_2, v_3, v_4]
= \frac{\pi}{3}\Bigl(4\pi - \sum_{i<j}\varphi_{ij} - \sum_\ell \arccos\frac{d_\ell}{R}\Bigr),
\]
where $\varphi_{ij}$ are the Euclidean dihedral angles of the tetrahedron
$T = \mathrm{conv}(v_1,\ldots,v_4)$, $R$ is its circumradius, and $d_\ell$
is the signed distance from the circumcenter to face $\ell$.

The \emph{boundary $3$-measure} (surface area) is the sum of hyperbolic
3-volumes of the boundary facets: each facet is a tetrahedron on $S^3$,
hence an ideal tetrahedron in a totally geodesic $\mathbb{H}^3 \subset \mathbb{H}^4$.
Its $\mathbb{H}^3$-volume is computed via the Lobachevsky formula
$\Lambda(\alpha) + \Lambda(\beta) + \Lambda(\gamma)$, where $\alpha,\beta,\gamma$
are the angles of a triangle obtained by stereographic projection from one
facet vertex.

\begin{remark}
In $\mathbb{H}^3$, the boundary of an ideal polyhedron consists of ideal
polygons, each with $\mathbb{H}^2$-area $(k-2)\pi$ for a $k$-gon---a
topologically determined quantity. In $\mathbb{H}^4$, boundary facets are ideal
tetrahedra with continuously varying $\mathbb{H}^3$-volumes. Thus the surface
area optimization problem is genuinely geometric only in dimension $d \geq 4$.
\end{remark}

\subsection{Combinatorial Invariance of Optimal Polytopes}

We numerically optimized configurations of $N$ points on $S^{d-1}$
($d = 3, 4$) to maximize three functionals:
\begin{enumerate}[label=(\alph*)]
\item \textbf{Volume:} Euclidean volume of $\mathrm{conv}(v_1,\ldots,v_N)$,
      or $\mathbb{H}^d$-volume of the ideal polytope.
\item \textbf{Surface area} ($d=4$ only): boundary $3$-measure.
\item \textbf{Isoperimetric ratio:} $I(P) = \mathrm{Vol}_d(P) / \mathrm{Area}_{d-1}(\partial P)$.
\end{enumerate}
For each $(N, \text{metric})$ pair, we performed multi-start optimization
in both the Euclidean and hyperbolic settings,
then compared the $f$-vectors of the optimal convex hulls.
The Euclidean phase uses the CMA-ES (Covariance Matrix Adaptation
Evolution Strategy) algorithm~\cite{HansenOstermeier2001},
which adapts a multivariate Gaussian to the promising region of
parameter space.  This dramatically outperforms random restarts:
in $d = 4$, CMA-ES finds Euclidean-optimal configurations that
are within $0.0$--$0.4\%$ of the hyperbolic optimum after a single
derivative-free Powell polish.  The hyperbolic phase uses either
Powell's method ($d = 4$, where volume evaluation via the
Haagerup--Munkholm formula takes $\sim$2\,ms) or torch LBFGS
with exact gradients through a custom $\mathrm{Li}_2$
autograd function ($d = 6$, where the Hopf--Murakami formula
requires dilogarithm evaluations).

\begin{table}[ht]
\centering
\caption{Combinatorial type comparison: Euclidean vs.\ Hyperbolic volume maximization in $d = 4$ (points on $S^3$).
CMA-ES Euclidean warm start with Powell hyperbolic polish.
Exact $f$-vector match holds for $N \leq 40$; small discrepancies
($\leq 3$ facets) appear at larger~$N$.}
\label{tab:comb_d4}
\smallskip
\begin{tabular}{r l l c r r}
\toprule
$N$ & Eucl.\ $f$-vector & Hyp.\ $f$-vector & Match & Eucl.\ vol & Hyp.\ vol \\
\midrule
5  & $(5,10,10,5)$     & $(5,10,10,5)$     & Yes & 0.146 & 0.269 \\
6  & $(6,15,18,9)$     & $(6,15,18,9)$     & Yes & 0.281 & 0.588 \\
7  & $(7,19,24,12)$    & $(7,19,24,12)$    & Yes & 0.433 & 1.024 \\
8  & $(8,24,32,16)$    & $(8,24,32,16)$    & Yes & 0.667 & 1.771 \\
9  & $(9,29,40,20)$    & $(9,29,40,20)$    & Yes & 0.793 & 2.314 \\
10 & $(10,35,50,25)$   & $(10,35,50,25)$   & Yes & 0.908 & 2.894 \\
12 & $(12,46,68,34)$   & $(12,46,68,34)$   & Yes & 1.158 & 4.324 \\
16 & $(16,69,104,52)$  & $(16,69,104,52)$  & Yes & 1.598 & 7.794 \\
20 & $(20,92,144,72)$  & $(20,92,144,72)$  & Yes & 1.958 & 11.936 \\
24 & $(24,117,186,93)$ & $(24,117,186,93)$ & Yes & 2.175 & 15.568 \\
30 & $(30,153,248,124)$& $(30,153,248,124)$& Yes & 2.479 & 22.068 \\
\midrule
40 & $(40,220,360,180)$& $(40,220,360,180)$& Yes & 2.813 & 32.936 \\
50 & $(50,282,464,232)$& $(50,281,462,231)$& \textbf{No} & 3.069 & 45.014 \\
60 & $(60,347,574,287)$& $(60,346,572,286)$& \textbf{No} & 3.243 & 56.711 \\
80 & $(80,478,796,398)$& $(80,478,796,398)$& Yes & 3.506 & 81.416 \\
100& $(100,613,1026,513)$& $(100,610,1020,510)$& \textbf{No} & 3.685 & 107.411 \\
\bottomrule
\end{tabular}
\end{table}

\begin{table}[ht]
\centering
\caption{Same comparison in $d=3$ (points on $S^2$). All $f$-vectors match.}
\label{tab:comb_d3}
\smallskip
\begin{tabular}{r l c r r}
\toprule
$N$ & $f$-vector & Match & Eucl.\ vol & Hyp.\ vol \\
\midrule
4  & $(4,6,4)$       & Yes & 0.513 & 1.015 \\
6  & $(6,12,8)$      & Yes & 1.333 & 3.664 \\
8  & $(8,18,12)$     & Yes & 1.816 & 6.489 \\
12 & $(12,30,20)$    & Yes & 2.536 & 13.529 \\
20 & $(20,54,36)$    & Yes & 3.119 & 27.673 \\
30 & $(30,84,56)$    & Yes & 3.450 & 46.231 \\
50 & $(50,144,96)$   & Yes & 3.733 & 83.990 \\
\bottomrule
\end{tabular}
\end{table}

The key empirical finding is:

\begin{conjecture}[Combinatorial invariance --- strong form]
\label{conj:comb_invariance}
For $N$ points on $S^{d-1}$, the $f$-vector of the volume-maximizing
convex polytope is the same whether volume is measured in the Euclidean
sense (volume of the convex hull) or in the hyperbolic sense (volume of the
ideal polytope in $\mathbb{H}^d$).
\end{conjecture}

This holds exactly in all tested cases with $N \lesssim 5d$:
$d=3$ ($N \leq 50$), $d=4$ ($N \leq 40$), $d=5$ ($N \leq 20$),
and $d=6$ ($N \leq 30$).
However, extending to larger~$N$ reveals a more nuanced picture.
At $N = 50$ in $d = 4$ and $N = 40$ in $d = 6$, the hyperbolic optimizer
finds a slightly different combinatorial type---with $1$--$3$ fewer facets
than the Euclidean optimum.  The discrepancies are small
($\leq 0.6\%$ of the facet count) and non-monotone: the $f$-vectors
match again at $N = 80$ in $d = 4$ and at $N = 50$ in $d = 6$.
In every case of divergence, the hyperbolic optimizer has \emph{fewer}
facets, paralleling the surface-area vs.\ volume divergence
(Section~\ref{sec:cross_metric}).
This suggests a weaker form that may hold universally:

\begin{conjecture}[Combinatorial invariance --- weak form]
\label{conj:comb_weak}
For $N$ points on $S^{d-1}$, the $f$-vectors of the Euclidean and
hyperbolic volume maximizers satisfy
$|f_i^{\mathrm{Eucl}} - f_i^{\mathrm{Hyp}}| = O(1)$
as $N \to \infty$, with $f_i^{\mathrm{Hyp}} \leq f_i^{\mathrm{Eucl}}$
for all~$i$.
\end{conjecture}

The CMA-ES Euclidean warm start is within $0.0$--$0.4\%$ of the
hyperbolic optimum in every case, confirming that the two objectives
share essentially the same landscape near their maxima.

\subsection{Cross-Metric Divergence}
\label{sec:cross_metric}

While the combinatorial type is invariant under change of geometry
(Euclidean vs.\ hyperbolic), it is \emph{not} invariant under change
of objective function. In $d=4$, the surface-area-maximizing polytope
can have a different $f$-vector from the volume-maximizing one:

\begin{center}
\begin{tabular}{r l l}
\toprule
$N$ & Volume-optimal $f$-vector & Area-optimal $f$-vector \\
\midrule
6  & $(6,15,18,9)$   & $(6,14,16,8)$ \\
16 & $(16,69,106,53)$ & $(16,68,104,52)$ \\
30 & $(30,154,248,124)$ & $(30,153,246,123)$ \\
\bottomrule
\end{tabular}
\end{center}

In each divergent case, the area-optimal polytope has fewer edges and facets.
Geometrically, surface area is maximized when boundary facets approach the
regular ideal tetrahedron volume $v_3 \approx 1.0149$, which favors fewer,
more regular facets; volume maximization instead favors packing more
simplicial cells into the interior.

\subsection{Correlation Between Euclidean and Hyperbolic Volumes}

The combinatorial invariance suggests a deeper relationship between
the two volume functionals. To quantify this, we computed both
$V_E$ (Euclidean hull volume) and $V_H$ (hyperbolic ideal polytope volume)
for 200 independent local optima at each $N$, obtained by single-phase
L-BFGS-B from random starting points.

\begin{center}
\begin{tabular}{r r r r}
\toprule
$N$ & Correlation & Log-log elasticity & $V_H/V_E$ (mean $\pm$ std) \\
\midrule
10 & 0.9985 & 1.21 & $3.17 \pm 0.02$ \\
20 & 0.9986 & 1.66 & $5.97 \pm 0.06$ \\
30 & 0.9956 & 1.97 & $8.78 \pm 0.05$ \\
\bottomrule
\end{tabular}
\end{center}

The Pearson correlation exceeds 0.995 in all cases: in the space of
near-optimal configurations, the Euclidean and hyperbolic volumes are
almost perfectly correlated. The log-log elasticity
$\partial \log V_H / \partial \log V_E$ increases with $N$,
approaching 2 at $N = 30$. This means a 1\% improvement in
Euclidean volume yields approximately a 2\% improvement in
hyperbolic volume near the optimum---consistent with empirical
observations from independent optimization.

This near-proportionality is surprising: $V_E$ is a polynomial in
the vertex coordinates, while $V_H$ involves the transcendental
Haagerup--Munkholm formula. The correlation suggests that both
functionals are effectively measuring the same underlying geometric
quantity---roughly, how ``uniformly spread'' the points are on
$S^{d-1}$.

\subsection{Non-Isolated Maxima: The Volume Plateau}

A striking observation emerges at $N = 9$ (equivalently, 20 simplicial facets
in $d = 4$). We performed two independent 500-restart optimizations:
\begin{enumerate}[label=(\roman*)]
\item Maximize \emph{Euclidean} hull volume $V_E$ over 9 points on $S^3$.
\item Maximize \emph{hyperbolic} ideal polytope volume $V_H$ over 9 points on $S^3$.
\end{enumerate}
Both converge to the same hyperbolic volume
$V_H = 2.313851299873$ (to 13 digits), yet the two
optimal configurations are \emph{metrically distinct}: their sorted
pairwise chordal distance vectors differ by up to $0.25$.
A third configuration, found independently by GPT-o3
(using a different optimization method), achieves
$V_H = 2.313850783625$---the same to 6 digits but a demonstrably
different polytope (different Gram matrix).

The Euclidean maximizer turns out to be a classical algebraic object:
the \emph{free sum} $C_5 \oplus I \oplus I$, the convex hull of a
regular pentagon in the $(x_1,x_2)$-plane and two orthogonal diameters
$\pm e_3$, $\pm e_4$.  Its Euclidean volume and surface area are
\[
V_E^* = \tfrac{5}{24}\sqrt{10 + 2\sqrt{5}}
      \approx 0.7925,
\qquad
A_E^* = \tfrac{5}{3}\sqrt{15 - \sqrt{5}}
      \approx 5.9544.
\]
The free sum simultaneously maximizes \emph{both} the Euclidean
4-volume and the boundary 3-measure among all 9-vertex polytopes
inscribed in $S^3$ (500 independent restarts of each objective
all converge to these values).  It has only four distinct pairwise
chordal distances, with $\sqrt{2}$ appearing 24 out of 36 times.

More generally, for any composite number $n = pq$ of facets,
the free sum $C_p \oplus C_q$ (a $p$-gon and $q$-gon in orthogonal
planes, with all vertices on $S^3$) gives a simplicial 4-polytope
with exactly $pq$ tetrahedral facets and $p + q$ vertices.
For prime $n$, no such product construction is available, and the
optimal configuration may lack algebraic structure.

A systematic sweep over all facet counts $9 \leq n \leq 37$ reveals a
richer picture.  The free sum $C_p \oplus C_q$ requires both $p, q \geq 3$,
so it is available precisely when $n$ admits a factorization $n = pq$ with
$p, q \geq 3$---excluding primes, twice-primes ($2p$), and $n < 9$.
For each $n$, we also performed 200-restart penalized L-BFGS-B optimization
over all inscribed simplicial 4-polytopes with exactly $n$ facets.
Key findings:

\begin{enumerate}[label=(\roman*)]
\item \textbf{Balanced products dominate:} At fixed facet count, the free
      sum with the most balanced factorization gives the largest volume.
      For example, $C_4 \oplus C_5$ ($n = 20$, vol $\approx 0.793$) far exceeds
      $C_3 \oplus C_7$ ($n = 21$, vol $\approx 0.592$).
\item \textbf{Unbalanced products are suboptimal:} For $n = 21$ and $n = 27$
      ($= 9 \times 3$), the unconstrained optimizer beats the free sum
      by 28\% and 39\% respectively, finding asymmetric configurations
      with higher volume.
\item \textbf{Volume grows primarily with vertex count:} At fixed vertex
      count $N$, volume increases monotonically with facet count.
      Products pack more facets per vertex than generic configurations
      (e.g., $C_5 \oplus C_5$ achieves 25 facets with only 10 vertices),
      which is the primary source of their advantage.
\item \textbf{Primes do not intrinsically trail:} At fixed $N$,
      prime facet counts are only slightly below neighboring composites
      (e.g., at $N = 9$: $f = 23$ (prime) gives vol $\approx 0.743$
      vs.\ $f = 22$ ($= 2 \times 11$) at $0.757$ and
      $f = 21$ ($= 3 \times 7$) at $0.758$).
      The apparent deficit of primes in the overall ranking is an artifact
      of the facet--vertex relationship, not number theory.
\end{enumerate}

\subsection{Higher Dimensions}

Extending the Euclidean cross-metric comparison (volume vs.\ surface area)
to $d = 5$ and $d = 6$, we find the same qualitative pattern:

\begin{center}
\begin{tabular}{r r l l c}
\toprule
$d$ & $N$ & Vol-optimal $f$-vector & Area-optimal $f$-vector & Match \\
\midrule
5 & 6  & $(6,15,6)$ & $(6,15,6)$ & Yes \\
5 & 7  & $(7,21,12)$ & $(7,21,12)$ & Yes \\
5 & 8  & $(8,27,18)$ & $(8,26,16)$ & \textbf{No} \\
5 & 9  & $(9,33,24)$ & $(9,33,24)$ & Yes \\
5 & 10 & $(10,40,32)$ & $(10,40,32)$ & Yes \\
5 & 15 & $(15,80,82)$ & $(15,80,82)$ & Yes \\
\midrule
6 & 7  & $(7,21,7)$ & $(7,21,7)$ & Yes \\
6 & 8  & $(8,28,16)$ & $(8,28,15)$ & \textbf{No} \\
6 & 9  & $(9,36,27)$ & $(9,36,27)$ & Yes \\
6 & 10 & $(10,43,36)$ & $(10,43,36)$ & Yes \\
6 & 12 & $(12,60,64)$ & $(12,60,64)$ & Yes \\
6 & 18 & $(18,123,196)$ & $(18,123,198)$ & \textbf{No} \\
\bottomrule
\end{tabular}
\end{center}

In every dimension $d \geq 4$, volume and surface area optimizers
agree for most values of $N$, but diverge at specific small $N$
(around $d+2$ to $d+3$). The direction of divergence is universal:
\emph{the area optimizer always finds a polytope with fewer or equal
facets compared to the volume optimizer}. In $d = 3$, no divergence
was observed for $N \leq 50$.

\subsection{Ideal Polytopes in $\mathbb{H}^5$}

Computing the volume of an ideal simplex in $\mathbb{H}^5$ is more
challenging than in even dimensions: there is no closed-form analogue
of the Hopf reduction.  Instead, the volume must be evaluated via the
Schl\"afli path integral.  However, for \emph{optimization}, the volume
itself is not needed: the Schl\"afli differential formula
\[
dV_5 = -\tfrac{1}{4}\sum_{F} V_3(F)\,d\theta_F
\]
gives the exact gradient in closed form, since each codimension-2 face~$F$
of an ideal 5-simplex is an ideal tetrahedron in a totally geodesic
$\mathbb{H}^3$, with volume given by the Milnor--Lobachevsky formula
$V_3 = \Lambda(\alpha) + \Lambda(\beta) + \Lambda(\gamma)$.
The dihedral angles~$\theta_F$ come from the angle Gram matrix.
Thus the optimization loop is entirely quadrature-free; the
Schl\"afli path integral is evaluated only once at the final
configuration to report the volume.

\begin{table}[ht]
\centering
\caption{Combinatorial invariance: Euclidean vs.\ ideal $\mathbb{H}^5$
volume maximization (points on $S^4$).  All $f$-vectors match,
and the hyperbolic polish improvement is exactly zero in every case.}
\label{tab:comb_H5}
\smallskip
\begin{tabular}{r l c r r}
\toprule
$N$ & $f$-vector & Match & Eucl.\ vol & $\mathbb{H}^5$ vol \\
\midrule
7  & $(7,21,12)$      & Yes & 0.06667 & 0.13123 \\
8  & $(8,27,18)$      & Yes & 0.11250 & 0.24414 \\
9  & $(9,33,24)$      & Yes & 0.17321 & 0.41394 \\
10 & $(10,40,32)$     & Yes & 0.26667 & 0.69954 \\
12 & $(12,55,50)$     & Yes & 0.37688 & 1.13886 \\
14 & $(14,71,70)$     & Yes & 0.48199 & 1.64660 \\
16 & $(16,89,94)$     & Yes & 0.58601 & 2.23211 \\
20 & $(20,125,142)$   & Yes & 0.80565 & 3.75231 \\
\bottomrule
\end{tabular}
\end{table}

In contrast to $d = 4$ and $d = 6$, where the hyperbolic polish
improves the volume by $0.03$--$0.4\%$, in $d = 5$ the CMA-ES
Euclidean optimizer already finds the hyperbolic optimum to full
precision.  This suggests that the Euclidean and hyperbolic volume
landscapes on $S^4$ are even more tightly aligned than in even
dimensions.

\paragraph{Reproducibility note on Table~\ref{tab:comb_H5}.}
Vertex coordinates for the $\mathbb{H}^5$-optimal polytopes in the
table were not retained alongside the volumes at the time of the
original runs.  We subsequently recomputed $\mathbb{H}^5$ volumes
via the Schl\"afli path integral from the stored
Euclidean-optimal coordinates in
\texttt{optimal\_d5\_d6\_symmetry\_analysis.json}.  The recomputed
values match Table~\ref{tab:comb_H5} to $\sim 0.1\%$ (paper
print precision) at $N \in \{7, 8, 9, 10, 16\}$, confirming that
those stored Euclidean-optima are also $\mathbb{H}^5$-optima.
At $N \in \{12, 14, 20\}$ the recomputed volumes fall short of
the table by $1$--$4\%$, indicating that the retained coordinates
at those $N$ represent local rather than global maxima of the
Euclidean objective; the values in the table, produced at the
original-run budget, are the tighter records and are treated as
authoritative here.  The symmetry structure of Table~\ref{tab:comb_H5}
(join-of-cyclic-polygons at $N \leq 11$, trivial join structure at
$N \geq 12$) is consistent with this: the high-symmetry
configurations were easier for the optimizer to locate
reproducibly.

\subsection{Ideal Polytopes in $\mathbb{H}^6$}

We extended the combinatorial invariance test to $d = 6$ by computing
ideal $\mathbb{H}^6$ polytope volumes directly.  For ideal simplices,
the vertex-link terms vanish ($w_0 = 0$), and the Hopf reduction gives
\[
\operatorname{Vol}_6(\Delta_{\mathrm{ideal}})
= \frac{8\pi^3}{15}\left(\frac{17}{4} - w_4 + \frac{1}{2}w_2\right),
\]
where $w_4$ involves only dihedral angles and $w_2$ involves only
$\mathbb{S}^3$-tetrahedral link volumes computable by Murakami's
dilogarithm formula.  A float64 implementation of Murakami's formula
achieves a 108$\times$ speedup over arbitrary-precision evaluation,
making optimization feasible.  Full details are in the companion
note~\cite{H56note}.

\begin{table}[ht]
\centering
\caption{Combinatorial invariance: Euclidean vs.\ ideal $\mathbb{H}^6$
volume maximization (points on $S^5$).  Exact match holds for $N \leq 30$;
at $N = 40$ the hyperbolic optimizer finds $2$ fewer facets.}
\label{tab:comb_H6}
\smallskip
\begin{tabular}{r l l c r r}
\toprule
$N$ & Eucl.\ $f$-vector & Hyp.\ $f$-vector & Match & Eucl.\ vol & $\mathbb{H}^6$ vol \\
\midrule
7  & $(7,21,7)$       & $(7,21,7)$       & Yes & 0.00584 & 0.01023 \\
8  & $(8,28,16)$      & $(8,28,16)$      & Yes & 0.01317 & 0.02504 \\
9  & $(9,36,27)$      & $(9,36,27)$      & Yes & 0.02436 & 0.05022 \\
10 & $(10,43,36)$     & $(10,43,36)$     & Yes & 0.03750 & 0.08390 \\
12 & $(12,60,64)$     & $(12,60,64)$     & Yes & 0.08889 & 0.23284 \\
14 & $(14,79,100)$    & $(14,79,100)$    & Yes & 0.12563 & 0.36922 \\
16 & $(16,100,145)$   & $(16,100,145)$   & Yes & 0.16588 & 0.54078 \\
20 & $(20,146,251)$   & $(20,146,251)$   & Yes & 0.24508 & 0.95399 \\
24 & $(24,199,390)$   & $(24,199,390)$   & Yes & 0.32510 & 1.47180 \\
30 & $(30,284,620)$   & $(30,284,620)$   & Yes & 0.45559 & 2.51089 \\
\midrule
40 & $(40,435,1049)$  & $(40,435,1047)$  & \textbf{No} & 0.64425 & 4.62319 \\
50 & $(50,608,1576)$  & $(50,608,1576)$  & Yes & 0.81844 & 7.20649 \\
\bottomrule
\end{tabular}
\end{table}

The strong form of the combinatorial invariance conjecture
(Conjecture~\ref{conj:comb_invariance}) holds exactly in
$d = 3$ ($N \leq 50$), $d = 4$ ($N \leq 40$), $d = 5$ ($N \leq 20$),
and $d = 6$ ($N \leq 30$).
At larger~$N$, small deviations appear (see Tables~\ref{tab:comb_d4}
and~\ref{tab:comb_H6}), consistent with the weak form
(Conjecture~\ref{conj:comb_weak}).

\subsection{Identification of Optimal Polytopes: Joins of Regular Polygons}
\label{sec:joins}

Searching polytope databases and comparing with known constructions,
we find that the volume-optimal polytopes for small~$N$ are
\emph{joins (free sums) of regular polygons in mutually orthogonal $2$-planes}.

In $d = 2k$, the join $C_{p_1} \oplus C_{p_2} \oplus \cdots \oplus C_{p_k}$
of $k$ regular polygons in $k$ orthogonal $2$-planes in $\mathbb{R}^{2k}$
gives a simplicial polytope with $p_1 + \cdots + p_k$ vertices on
$S^{2k-1}$ and $p_1 \cdots p_k$ tetrahedral facets. The following table
identifies each small-$N$ optimal polytope:

\begin{center}
\begin{tabular}{r l l l}
\toprule
$N$ & $d=4$ optimal & $d=6$ optimal & Structure \\
\midrule
$d+1$ & $(5,10,10,5)$ & $(7,21,7)$ & Simplex $\Delta_d$ \\
$d+2$ & $(6,15,18,9)$ & $(8,28,16)$ & $\Delta_{d/2} * \Delta_{d/2}$ \\
$d+3$ & $(7,19,24,12)$ & $(9,36,27)$ & $C_3^{d/2}$ (join of $d/2$ triangles) \\
$d+4$ & $(8,24,32,16)$ & $(10,43,36)$ & Cross-polytope at $d{=}4$; $C_3^2{\oplus}C_4$ at $d{=}6$\\
$d+5$ & $(9,29,40,20)$ & --- & $C_4{\oplus}C_5$ \\
$d+6$ & $(10,35,50,25)$ & $(12,60,64)$ & $C_5^2$ at $d{=}4$; $C_4^3$ (cross-polytope) at $d{=}6$ \\
\bottomrule
\end{tabular}
\end{center}

For $d = 4$, the optimal polytope matches the best balanced
free sum $C_p \oplus C_q$ for all $N \leq 10$, but not for $N \geq 12$.
For $d = 6$, the optimal matches the best balanced triple
$C_p \oplus C_q \oplus C_r$ for all $N \leq 14$, but not for $N = 16$.
In both cases the threshold is approximately $N \approx 2d + 2$, beyond
which the optimizer discovers asymmetric configurations that beat
the rigid product structure.

Notable special cases:
\begin{itemize}
\item $N = 2d$: the \emph{$d$-dimensional cross-polytope}
      ($C_4 \oplus C_4$ in $d = 4$; $C_4 \oplus C_4 \oplus C_4$ in $d = 6$).
\item $N = d + 2$: the join of two simplices of half the dimension
      ($\Delta_2 * \Delta_2$ in $d = 4$; $\Delta_3 * \Delta_3$ in $d = 6$).
\end{itemize}

The volume-optimal polytopes are \emph{not} cyclic polytopes, and
are in fact far from neighborly: at $N = 30$ in $d = 4$, the optimal
polytope has only 123 facets compared to 756 for the cyclic polytope
$C(30,4)$. Maximizing facet count and maximizing volume are almost
opposite goals.

The hyperbolic-optimal configuration, by contrast, has no transposition
symmetries and a completely different distance spectrum.
Yet both achieve the same hyperbolic volume $V_H = 2.31385$ to
13 significant digits. A third configuration, found independently by
GPT-o3 using a different optimization method, achieves
$V_H = 2.31385078$---the same to 6 digits but again a
demonstrably different polytope (different Gram matrix).

The Hessian of $V_H$ at the hyperbolic optimum
has two near-zero eigenvalues ($\sim 10^{-3}$), indicating that
the maximum is achieved on (at least) a 2-dimensional manifold
in the 16-dimensional configuration space
$\mathrm{Conf}_9(S^3) / \mathrm{O}(4)$.

This suggests a phenomenon that is invisible to
$f$-vector comparisons: not only do different metrics find the same
combinatorial type, but the hyperbolic volume maximum is achieved on
a \emph{locus} that contains configurations with very different
Euclidean properties.  The free sum $C_5 \oplus I \oplus I$---the
unique maximizer of $V_E$---and the asymmetric hyperbolic maximizer
are different points on this locus, connected by a flat ridge along
which $V_H$ is constant (to the precision of our computation)
but $V_E$ varies by over 10\%.

\subsection{Symmetry Analysis via Douady--Earle Centering}

Beyond $f$-vector comparisons, we can determine the full geometric
automorphism group of each optimal polytope by normalizing the
Douady--Earle conformal barycenter to the origin of the ball model.
The residual symmetries are then elements of $\mathrm{O}(d)$, and
the automorphism group can be computed exactly using
\textsc{nauty}~\cite{McKayPiperno} on the edge-colored complete graph
whose colors encode the chordal distance classes.

\begin{table}[ht]
\centering
\caption{Symmetry groups of volume-optimal inscribed polytopes in $d = 4$.
For small~$N$, the join structure produces highly symmetric polytopes;
symmetry breaks sharply at $N \approx 11$ and vanishes beyond $N = 20$.}
\label{tab:symmetry_d4}
\smallskip
\begin{tabular}{r r r l l}
\toprule
$N$ & Facets & $|\mathrm{Aut}|$ & Group & Orbits \\
\midrule
5  & 5  & 120     & $S_5$            & $[5]$    \\
6  & 9  & 720     & $S_6$            & $[6]$    \\
7  & 12 & 48      & $D_4 \times S_3$ & $[4,3]$  \\
8  & 16 & 384     & $\beta_4$ (hyperoctahedral) & $[8]$    \\
9  & 20 & 80      & $D_4 \times D_5$ & $[5,4]$  \\
10 & 25 & 200     & $(C_5 \times C_5) \rtimes D_4$ & $[10]$   \\
\midrule
11 & 29 & 8       & $D_4$            & $[4,4,2,1]$ \\
12 & 34 & 8       & $D_4$            & $[8,4]$  \\
\textbf{13} & \textbf{39} & \textbf{52} & $\mathbf{C_{13} \rtimes C_4}$ & $\mathbf{[13]}$ \\
14 & 43 & 2       & $C_2$            & $[2^6,1^2]$ \\
15 & 48 & 72      & $(S_3 \times S_3) \rtimes C_2$ & $[6,6,2,1]$ \\
16 & 53 & 2       & $C_2$            & 9 orbits \\
20 & 72 & 16      & $C_2 \times D_4$ & 4 orbits \\
30 & 124& 1       & trivial          & all singletons \\
\bottomrule
\end{tabular}
\end{table}

The $N = 10$ case is noteworthy: the volume-optimal polytope is the
join $C_5 \oplus C_5$---two regular pentagons in orthogonal
$2$-planes---with vertex-transitive automorphism group
$(C_5 \times C_5) \rtimes D_4$ of order~200.  The $D_4$ factor swaps
the two pentagons and reflects within each; the $C_5 \times C_5$
factor acts by independent cyclic rotations.  This structure was
independently rediscovered by Google's Gemini~3.1~Pro model
(see~\S\ref{sec:llm_benchmark}), which identified the join
decomposition from the face lattice alone.

The most remarkable case is $N = 13$ in $d = 4$. The volume-maximizing
inscribed 4-polytope with 13 vertices on $S^3$ has a transitive
automorphism group of order~52 isomorphic to the Frobenius group
$C_{13} \rtimes C_4$, with only 3 distinct chordal distances.
This is the unique Frobenius group of order~52: the normal subgroup
$C_{13}$ acts as a cyclic permutation of all vertices, while the
complement $C_4$ acts as the subgroup of 4th powers in
$\mathrm{AGL}(1, \mathbb{F}_{13})$.
The polytope has 39 tetrahedral facets ($= 3 \times 13$).

This vertex-transitive polytope does not appear in any catalog of
uniform, scaliform, or swirlchoron polytopes that we could locate.
The classification of vertex-transitive convex 4-polytopes is an
open problem (unlike in 3D, where they are all deformations of
Archimedean solids).

\subsubsection{The Frobenius construction}

The $N = 13$ polytope belongs to an infinite family.  For each prime
$p \equiv 1 \pmod{4}$, let $k$ be a primitive 4th root of unity
modulo~$p$ (so $k^2 \equiv -1$).  The vertices
\[
  v_j = \frac{1}{\sqrt{2}}\Bigl(
    \cos\frac{2\pi j}{p},\;
    \sin\frac{2\pi j}{p},\;
    \cos\frac{2\pi kj}{p},\;
    \sin\frac{2\pi kj}{p}
  \Bigr), \qquad j = 0,\ldots,p-1,
\]
form a vertex-transitive polytope on $S^3$ with automorphism
group $\mathbb{Z}_p \rtimes \mathbb{Z}_4$ and exactly $(p-1)/4$ chordal distance
classes~\cite{FrobeniusNote}.  The construction arises from a
faithful four-dimensional real representation induced by a single orbit
of $\mathbb{Z}_4$ acting on the characters of~$\mathbb{Z}_p$.

\begin{center}
\begin{tabular}{r r r r}
\toprule
$p$ & $|{\mathrm{Aut}}|$ & Distance classes & Volume status \\
\midrule
5  & 20  & 1 & global max (= simplex $\Delta_4$) \\
13 & 52  & 3 & global max \\
17 & 68  & 4 & local max (not global) \\
29 & 116 & 7 & critical point \\
37 & 148 & 9 & critical point \\
\bottomrule
\end{tabular}
\end{center}

Every Frobenius configuration is a critical point of the hyperbolic volume
functional (the gradient vanishes by equivariance: the $\mathbb{Z}_4$ stabilizer
of the base vertex leaves no free parameter after normalization to $S^3$).
Hessian analysis shows that $p = 13$ and $p = 17$ are both
\emph{non-degenerate local maxima}: the constrained Hessian on
the tangent space to $(S^3)^p / \mathrm{O}(4)$ is negative-definite,
with exactly $\dim \mathrm{SO}(4) = 6$ zero eigenvalues.  For $p = 17$
the Frobenius polytope achieves $85\%$ of the global maximum volume,
which is attained by an asymmetric configuration with $|\mathrm{Aut}| = 1$.

\subsubsection{Volume ridges}

The Hessian analysis also reveals that the \emph{join} configurations
$C_p \oplus C_q$ have a degenerate local-maximum structure: the
hyperbolic volume Hessian at these configurations has
$\dim \mathrm{SO}(d) + r$ zero eigenvalues for some $r > 0$,
indicating a volume-constant ridge of dimension~$r$.  In $d = 4$,
every join has $r = 4$ extra zero eigenvalues; in $d = 6$, all
tested joins have $r = 6$.  Thus the hyperbolic volume maximum for
join-type polytopes is achieved on a positive-dimensional manifold
in the configuration space $\mathrm{Conf}_N(S^{d-1}) / \mathrm{O}(d)$,
containing configurations with very different Euclidean properties.
The Frobenius polytopes, by contrast, are isolated (non-degenerate)
local maxima.

The pattern is universal across dimensions $d = 4, 5, 6$:
\begin{enumerate}[label=(\roman*)]
\item \textbf{Join regime} ($N \leq \lfloor 2.5 d \rfloor$):
      the optimum is a join of regular polygons with large
      automorphism group ($|\mathrm{Aut}| \geq 100$).
      The cross-polytope $\beta_d$ at $N = 2d$ is always the
      symmetry peak.
\item \textbf{Sharp transition} (1--2 values of $N$):
      symmetry drops by 1--3 orders of magnitude.
\item \textbf{Asymmetric regime} ($N \geq 4d$--$5d$):
      $|\mathrm{Aut}| = 1$, all vertices in distinct orbits.
\end{enumerate}

\begin{table}[ht]
\centering
\caption{Symmetry groups of volume-optimal polytopes in $d = 5$
(points on $S^4$).  Joins $C_p \oplus C_q \oplus I$ dominate
for $N \leq 11$; symmetry breaks at $N = 12$.}
\label{tab:symmetry_d5}
\smallskip
\begin{tabular}{r r r l l r}
\toprule
$N$ & Facets & $|\mathrm{Aut}|$ & Structure & Orbits & \#dist \\
\midrule
6  & 6    & 720    & $\Delta_5$             & $[6]$    & 1 \\
8  & 18   & 144    & $C_3 \oplus C_3 \oplus I$ & $[6,2]$  & 3 \\
9  & 24   & 288    & $C_3 \oplus C_4 \oplus I$ & $[6,3]$  & 3 \\
10 & 32   & 3840   & $\beta_5$              & $[10]$   & 2 \\
11 & 40   & 480    & $C_4 \oplus C_5 \oplus I$ & $[6,5]$  & 4 \\
\midrule
12 & 48   & 64     & ?                      & $[4,4,4]$& 8 \\
13 & 58   & 16     & ---                    & $[4,4,2^2,1]$ & 20 \\
$\geq 15$ & --- & 1 & trivial & all singletons & --- \\
\bottomrule
\end{tabular}
\end{table}

\begin{table}[ht]
\centering
\caption{Symmetry groups of volume-optimal polytopes in $d = 6$
(points on $S^5$).  Triple joins $C_p \oplus C_q \oplus C_r$
dominate for $N \leq 15$; symmetry breaks at $N = 16$.}
\label{tab:symmetry_d6}
\smallskip
\begin{tabular}{r r r l l r}
\toprule
$N$ & Facets & $|\mathrm{Aut}|$ & Structure & Orbits & \#dist \\
\midrule
7  & 7     & 5040   & $\Delta_6$             & $[7]$    & 1 \\
8  & 16    & 1152   & $(S_4 \times S_4) \rtimes C_2$ & $[8]$ & 2 \\
9  & 27    & 1296   & $C_3^3$                & $[9]$    & 2 \\
10 & 36    & 576    & $C_3^2 \oplus C_4$     & $[6,4]$  & 3 \\
11 & 48    & 2304   & $C_3 \oplus C_4^2$     & $[8,3]$  & 3 \\
12 & 64    & 46080  & $\beta_6 = C_4^3$      & $[12]$   & 2 \\
13 & 80    & 3840   & $C_4^2 \oplus C_5$     & $[8,5]$  & 4 \\
14 & 100   & 1600   & $C_4 \oplus C_5^2$     & $[10,4]$ & 4 \\
15 & 125   & 6000   & $C_5^3$                & $[15]$   & 3 \\
\midrule
16 & 140   & 16     & ---                    & $[4,4,2^4]$ & 34 \\
$\geq 17$ & --- & $\leq 2$ & --- & --- & --- \\
\bottomrule
\end{tabular}
\end{table}

The classification of finite groups realizable as automorphism groups
of convex polytopes is now well understood:
every finite group arises as the full automorphism group of some
convex polytope~\cite{SchulteWilliams2015,Doignon2018}.
For \emph{vertex-transitive} polytopes, Babai~\cite{Babai1977}
showed that abelian groups other than elementary 2-groups, and
generalized dicyclic groups, cannot occur.
The Frobenius group $C_{13} \rtimes C_4$ at $N = 13$ in $d = 4$
passes Babai's test, confirming it is \emph{allowed} as a
vertex-transitive symmetry group---but to our knowledge it has
not previously been realized as a convex 4-polytope.

\subsection{Isoperimetric Ratio}

For the isoperimetric ratio $I(P) = \mathrm{Vol}_4(P) / \mathrm{Area}_3(\partial P)$,
the optimized values grow with $N$ but remain well below the Euclidean
ball ratio of $1/4$:

\begin{center}
\begin{tabular}{r r r}
\toprule
$N$ & Best $I(P)$ & Euclidean ball $I$ \\
\midrule
5  & 0.053 & \\
8  & 0.109 & \\
10 & 0.126 & \\
20 & 0.169 & \\
30 & 0.187 & 0.250 \\
\bottomrule
\end{tabular}
\end{center}

The hyperbolic ball of radius $r$ in $\mathbb{H}^4$ has
$\mathrm{Vol}_4 = (\pi/2)(\sinh 2r - 2r)$ and
$\mathrm{Area}_3 = 2\pi^2 \sinh^3 r$, giving
$I_{\mathrm{ball}}(r) \to 1/4$ as $r \to 0$ (Euclidean regime)
and $I_{\mathrm{ball}}(r) \sim e^{-r}/\pi \to 0$ as $r \to \infty$.
Whether the optimal $I(P)$ converges to a definite limit as
$N \to \infty$---and whether that limit is strictly less than
$1/4$---remains open.

% ============================================================================
\section{LLM Benchmarks: A Pointer}
\label{sec:llm_benchmark}
% ============================================================================

The objects studied here---volume-maximizing ideal polytopes, the
combinatorial structures they realize, and the associated face
lattices---turn out to provide a natural family of benchmarks for
large language models, with the property that an answer is a list of
vertex coordinates verifiable by a standalone checker rather than a
prose argument.  Three distinct benchmark families have been
designed and released alongside this paper, each probing a different
reasoning regime: inscribed reconstruction (face lattice
$\to$ coordinates), windowed volume (hit a target measure to within a
tolerance), and volume lower bound (give the sharpest bound you can
backed by explicit coordinates).  Their detailed specifications,
grader implementations, and frontier-model evaluations are reported
in the companion note~\cite{LLMNote}.  We restrict ourselves here to
two observations that bear directly on the mathematical content of
the present paper.

\subsection{The reconstruction problem and rigidity}

The reconstruction task asks: given the $f$-vector, edge list, and
facet list of a convex polytope inscribed in $S^{d-1}$, produce
$N$ explicit coordinates realizing it.  The inscribed constraint
is essential---without it, the problem reduces to combinatorial
realization of simplicial polytopes, a substantially easier
question.  When LLMs were tested on instances drawn from our
volume-optimal polytopes (large automorphism order, often joins) and
from random sphere configurations (generically trivial automorphism
group), GPT-5.4 Pro succeeded at instances up to $N = 30$ with
$|\mathrm{Aut}| = 1$ via numerical constraint satisfaction;
Gemini~3.1~Pro DeepThink succeeded on highly symmetric instances by
algebraic recognition (e.g., identifying the $N = 10$ optimal
polytope as $C_5 \oplus C_5$) but failed on asymmetric random
instances.  The relevance for the present paper is that the
\emph{rigidity} of the configuration---the dimension of the space
of inscribed coordinate configurations realizing a given
combinatorial type---is exactly the empirical difficulty axis the
LLMs see.  Volume-optimal polytopes are maximally rigid in this
sense; the hardest reconstruction instances should be large $N$
with simultaneously high volume and low automorphism order, of
which the $N = 16$ optimal polytope ($|\mathrm{Aut}| = 2$, $53$
facets, volume per vertex $0.49$) is a natural candidate.

\subsection{Volume lower bound and our records}

The volume-lower-bound benchmark asks the model to give the sharpest
LB it can on $V^*(d, N)$ and back the claim with explicit vertex
coordinates.  Our published records (Tables~\ref{tab:maximal-volumes},
\ref{tab:comb_H5}, etc.) are the reference; an LLM submission
classifies as \emph{strong} if it matches or beats them, \emph{pass}
if it is within $95\%$, and \emph{fail} otherwise.  An LLM that
beats our records produces new mathematical content; the
benchmark's anti-hallucination design ensures that this is a
falsifiable claim.  Hyperbolic $d = 5$ is excluded by design: as
discussed in the reproducibility note accompanying
Table~\ref{tab:comb_H5}, our stored $d = 5$ configurations at
$N \in \{12, 14, 20\}$ are known local rather than global maxima
of the Euclidean (and hence hyperbolic) objective in $d = 5$, so
grading against them would be unreliable in that dimension.

% ============================================================================
\section{Open Questions}
% ============================================================================

\begin{enumerate}
    \item \textbf{Proof of Rational Angles}: Can we prove that maximal volume implies rational dihedral angles?

    \item \textbf{Denominator Patterns}: What determines the denominator $q$ for a given $n$? Our data suggests connections to $n-2$ or $n-3$ in some cases.

    \item \textbf{Uniqueness}: Are maximal volume configurations unique up to symmetry?

    \item \textbf{Asymptotics}: What is the growth rate of $V_{\max}(n)$ as $n \to \infty$?

    \item \textbf{Beta Distribution}: Why does the normalized volume follow a Beta distribution? Is the limiting mean exactly $\ln 2$, or merely close to it? What geometric principle underlies the linear scaling of $\alpha_n$ and $\beta_n$ with $n$?

    \item \textbf{Combinatorial Invariance}: The strong form (Conjecture~\ref{conj:comb_invariance}) holds for $N \lesssim 5d$ but fails at larger~$N$, while the weak form (Conjecture~\ref{conj:comb_weak}) holds in all tested cases.  What determines the threshold?  Both the Euclidean and hyperbolic volume functionals factor through the Delaunay triangulation of the stereographic projection, which depends only on the M\"obius geometry of the point configuration.  The small discrepancies at large~$N$ ($\leq 3$ facets out of hundreds) suggest near-degenerate facets that resolve differently under the two metrics.  Is the weak conjecture true, and is the $O(1)$ bound tight?

    \item \textbf{Cross-Metric Divergence}: Characterize which functionals on $S^{d-1}$ give the same optimal combinatorial type and which do not. Our data suggests volume and isoperimetric ratio agree, while surface area can diverge.

    \item \textbf{Hyperbolic Isoperimetric Limit}: Does $\sup_P I(P)$ over all $N$-vertex ideal polytopes in $\mathbb{H}^4$ converge as $N \to \infty$, and if so, is the limit strictly less than $1/4$?

    \item \textbf{Spectral Realizability of Volume-Optimal Polytopes}.
    Let~$P$ be a simplicial 4-polytope inscribed in $S^3$ that
    maximizes hyperbolic volume among all inscribed polytopes with
    the same number of vertices, and let $L = D - A$ be the graph
    Laplacian of the 1-skeleton of~$P$.  Let
    $\varphi\colon i \mapsto (v_1(i), v_2(i), v_3(i), v_4(i)) / \|{\cdot}\|$
    be the map sending each vertex to the (normalized) vector of its
    first four nontrivial Laplacian eigenvector entries.

    \emph{Empirical observation:} $\varphi$ produces a polytope on~$S^{d-1}$
    combinatorially equivalent to~$P$ in dimensions $d = 3$ and $d = 4$,
    but \emph{not} in $d \geq 5$.  The results are as follows:
    \begin{center}
    \begin{tabular}{r l l l}
    \toprule
    $d$ & Success range & Failure onset & Success rate \\
    \midrule
    3 & $N = 6$ to $30$ & $N = 5$ (degenerate) & 25/25 \\
    4 & $N = 5$ to $30$ & never (in tested range) & 14/14 \\
    5 & $N = 10, 11$ only & $N = 12$ & 2/15 \\
    6 & $N = 12\text{--}15$ only & $N = 16$ & 4/18 \\
    \bottomrule
    \end{tabular}
    \end{center}
    In $d = 4$, this includes the $N = 30$ case with trivial
    automorphism group and 124 facets.  The method also fails for
    generic inscribable polytopes in $d = 4$: random inscribable
    simplicial 4-polytopes consistently produce the wrong
    combinatorial type under spectral embedding.

    \emph{Explanation.}  The dimensional threshold $d \leq 4$
    coincides exactly with the range of the
    Blind--Mani--Kalai theorem~\cite{Kalai1988}: for
    simplicial polytopes in $d \leq 4$, the 1-skeleton
    determines the entire face lattice.  In $d \geq 5$, the
    1-skeleton is no longer sufficient, and distinct
    simplicial polytopes can share the same graph.  Since the
    spectral embedding uses only the 1-skeleton (via the
    graph Laplacian), it contains all the combinatorial
    information needed for realization in $d \leq 4$, but
    is provably incomplete in $d \geq 5$.

    The remaining mystery is why the spectral embedding
    succeeds for volume-\emph{optimal} inscribed polytopes
    but fails for generic inscribable ones in the same
    dimension.  One possible connection: volume-optimal
    configurations have vertices ``maximally spread'' on
    $S^{d-1}$ (approximate spherical designs), and the
    Laplacian eigenvectors also spread vertices according to
    graph distance.  For configurations that are already
    optimally spread, these two notions of dispersion may
    coincide.  A proof would yield a fast, deterministic
    algorithm for inscribed realization of volume-optimal
    polytopes in dimensions $d \leq 4$.

    For $d \geq 5$, a natural candidate is the Hodge
    1-Laplacian $L_1 = B_1 B_1^T + B_2 B_2^T$ acting on
    edges~\cite{Forman2003}, where $B_1$ and $B_2$ are the
    boundary operators of the simplicial complex.  The
    up-Laplacian $L_1^{\mathrm{up}} = B_2 B_2^T$ encodes
    triangle co-incidence, carrying information beyond the
    1-skeleton.  Preliminary experiments show that the naive
    approach (project $L_1$ eigenvectors back to vertices
    by averaging over incident edges) does not improve over
    $L_0$ in $d \geq 5$; the projection step discards the
    higher-order structure.  A more principled approach---using
    the edge-level embedding to reconstruct the Gram matrix
    directly, or using sheaf Laplacians that preserve the
    fiber structure---may succeed.

    \item \textbf{Inscribability in Dimension 4}.
    In $d = 3$, Rivin's linear program characterizes which
    simplicial polytopes are inscribable in~$S^2$.  No analogue
    is known in $d = 4$.  Our spectral realization algorithm
    and the cone/Volkov method (optimizing dihedral angle deficits
    around codimension-2 faces of the cone from the origin)
    provide practical heuristics for constructing inscribed
    realizations when they exist.  Can these be developed into
    a complete decision procedure?  A neural-network classifier
    trained on combinatorial features of face lattices, using
    the realization algorithms as an oracle for generating
    training data, might reveal patterns pointing toward a
    combinatorial characterization.
\end{enumerate}

% ============================================================================
\section{Application: High-Girth Cubic Graphs}
\label{sec:cages}
% ============================================================================

The optimization framework developed for ideal polytope volumes
applies more broadly to combinatorial search problems.  As a proof
of concept, we apply a similar strategy to the \emph{cage problem}:
finding the smallest $3$-regular (cubic) graph with girth~$g$.

\subsection{Method}

We maximize girth over cubic graphs on $N$ vertices using a two-phase
approach inspired by the Euclidean--hyperbolic pipeline for polytope
volumes:
\begin{enumerate}[label=(\roman*)]
\item \textbf{Log-determinant warm start.}  Maximize
      $\log \tau(G) = \log \det(\tilde L)$, where $\tilde L$ is the
      reduced Laplacian and $\tau(G)$ is the number of spanning trees.
      This smooth proxy correlates with graph expansion, which in turn
      correlates with girth.  Optimization is by simulated annealing
      with edge swaps (which preserve $3$-regularity).
\item \textbf{Cycle-breaking polish.}  Find a shortest cycle by BFS,
      then attempt to eliminate it by swapping one of its edges with
      a random non-cycle edge.  Repeat.
\end{enumerate}

The log-determinant phase plays the same role as the Euclidean
warm start for polytopes: it navigates the search toward a
well-structured region, from which the discrete girth-improving
moves can push to higher values.  For incremental evaluation,
edge swaps are rank-4 updates to the Laplacian, and the resulting
determinant change can be computed in $O(N^2)$ via Cholesky
rank-1 updates (with \texttt{numba} JIT acceleration), compared
to $O(N^3)$ for a full eigendecomposition.

\subsection{Results}

\begin{table}[ht]
\centering
\caption{Best girth found for cubic graphs on $N$ vertices,
using log-determinant warm start + cycle breaking
(100 trials, total budget 600 edge swaps per trial).
Known cage values are shown for reference.}
\label{tab:cages}
\smallskip
\begin{tabular}{r c c l}
\toprule
$N$ & Best girth & Cage girth & Graph \\
\midrule
10 & 5 & 5 & Petersen (optimal) \\
14 & 6 & 6 & Heawood (optimal) \\
24 & 7 & 7 & McGee (optimal) \\
30 & 7 & 8 & --- \\
40 & 7 & --- & --- \\
50 & 7 & --- & --- \\
60 & 8 & --- & --- \\
\bottomrule
\end{tabular}
\end{table}

The method rediscovers three known cages exactly: the Petersen
graph ($g = 5$, $N = 10$), the Heawood graph ($g = 6$, $N = 14$),
and the McGee graph ($g = 7$, $N = 24$).  It reaches girth~8 at
$N = 60$, twice the $(3,8)$-cage size of~30 (the Tutte--Coxeter
graph).  Girth~7 is achieved consistently for all $N \geq 24$.

% ============================================================================
% References
% ============================================================================

\begin{thebibliography}{99}

\bibitem{Rivin1994}
I.~Rivin, ``Euclidean structures on simplicial surfaces and hyperbolic volume,'' \emph{Ann. of Math.} \textbf{139} (1994), 553--580.

\bibitem{RivinCombOpt}
I.~Rivin, ``Combinatorial optimization in geometry,'' \emph{Adv. in Appl. Math.} \textbf{31} (2003), 242--271.

\bibitem{RivinTriangulations}
I.~Rivin, ``Triangulations into groups,'' arXiv:math/0510613, 2005.

\bibitem{RivinDjinn}
I.~Rivin, ``I Built a `Djinn' That Fixes My Code While I Sleep,'' \emph{Impure Thoughts} (Substack), December 2025. \url{https://impurethoughts.substack.com/}

\bibitem{Milnor}
J.~Milnor, ``Hyperbolic geometry: The first 150 years,'' \emph{Bull. Amer. Math. Soc.} \textbf{6} (1982), 9--24.

\bibitem{H56note}
I.~Rivin, ``Notes on hyperbolic 5- and 6-simplex volumes: formulas, reductions, and a practical code path,'' companion note, 2026.

\bibitem{LLMNote}
I.~Rivin, ``Hyperbolic polytope volume as an LLM benchmark suite: three tasks of increasing freedom,'' companion note, 2026.

\bibitem{RLNote}
I.~Rivin, ``When meta-learned optimization fails to beat round-robin: a validation study on hyperbolic polytope volume,'' companion note, 2026.

\bibitem{RivinBasmajian}
I.~Rivin, ``On Basmajian's identities, and other stories,'' arXiv:1404.1583, 2014.

\bibitem{HansenOstermeier2001}
N.~Hansen and A.~Ostermeier, ``Completely derandomized self-adaptation in evolution strategies,'' \emph{Evol.\ Comput.} \textbf{9} (2001), 159--195.

\bibitem{McKayPiperno}
B.~D.~McKay and A.~Piperno, ``Practical graph isomorphism, {II},'' \emph{J. Symbolic Comput.} \textbf{60} (2014), 94--112.

\bibitem{SchulteWilliams2015}
E.~Schulte and G.~I.~Williams, ``Polytopes with preassigned automorphism groups,'' \emph{Discrete Comput.\ Geom.} \textbf{54} (2015), 444--458.

\bibitem{Doignon2018}
J.-P.~Doignon, ``Any finite group is the group of some binary, convex polytope,'' \emph{Discrete Comput.\ Geom.} \textbf{59} (2018), 451--460.

\bibitem{Babai1977}
L.~Babai, ``Symmetry groups of vertex-transitive polytopes,'' \emph{Geom.\ Dedicata} \textbf{6} (1977), 331--337.

\bibitem{DouadyEarle}
A.~Douady and C.~J.~Earle, ``Conformally natural extension of homeomorphisms of the circle,'' \emph{Acta Math.} \textbf{157} (1986), 23--48.

\bibitem{FrobeniusNote}
I.~Rivin, ``Frobenius polytopes on $S^3$,'' companion note, 2026.

\bibitem{Firsching2017}
M.~Firsching, ``Realizability and inscribability for simplicial polytopes via optimization,'' \emph{Math.\ Program.} \textbf{166} (2017), 273--295. arXiv:1710.00979.

\bibitem{Kalai1988}
G.~Kalai, ``A simple way to tell a simple polytope from its graph,'' \emph{J. Combin.\ Theory Ser.~A} \textbf{49} (1988), 381--383.

\bibitem{Forman2003}
R.~Forman, ``Bochner's method for cell complexes and combinatorial Ricci curvature,'' \emph{Discrete Comput.\ Geom.} \textbf{29} (2003), 323--374.

\end{thebibliography}

\end{document}
